%% [ĐÃ SỬA] Thêm NEW-04 (nhập lịch ca), cổng quyết định G1, quy ước chủ task, tải theo vai trò, ghi chú phiên bản hiện có; cập nhật tổng ngày công 68,5; thêm chỉ số độ chính xác ở M4.
\section{Kế hoạch dự án và liên kết công việc}

\subsection{Nền tảng dự án}
\begin{tabularx}{\textwidth}{|>{\raggedright\arraybackslash}p{0.30\textwidth}|X|}
\hline
Kho mã nguồn (DutyLog) & \href{https://github.com/Big-Data-Club/BDCHub---DutyLog}{Mở kho mã nguồn} \\ \hline
GitHub Organization & \href{https://github.com/Big-Data-Club}{Mở tổ chức GitHub} \\ \hline
GitHub Project & \href{https://github.com/orgs/Big-Data-Club/projects/3}{Mở GitHub Project} \\ \hline
GitHub Wiki & \href{https://github.com/Big-Data-Club/BDCHub---DutyLog/wiki}{Mở GitHub Wiki} \\ \hline
Kế hoạch kỹ thuật & Tệp \emph{DutyLog 2.0 Product Engineering Roadmap} \\ \hline
Biểu mẫu/bộ câu hỏi khảo sát gốc \newline (dữ liệu thô, cần ẩn email trước khi công bố) & \href{https://docs.google.com/spreadsheets/d/1R5b7StzAAhT9WhOII3t3W2ZHKKUwrysQPayB1g6tjMU/edit?gid=1886321484\#gid=1886321484}{Mở bảng câu hỏi} \\ \hline
\end{tabularx}\par
% Khi công bố, đặt quyền chia sẻ của liên kết trên ở chế độ hạn chế (dữ liệu thô có email); bản ẩn danh và bộ câu hỏi nằm ở Phụ lục A.

\medskip
GitHub Project dùng luồng \textit{Backlog $\rightarrow$ To Do $\rightarrow$ In Progress $\rightarrow$ Review $\rightarrow$ Done}. Task được lấy từ roadmap kỹ thuật; chỉ các task phục vụ F1--F7 nằm ở To Do hoặc In Progress, các task còn lại giữ ở Backlog (mục 12.2).

\subsection{Các mốc công việc và phân công}

\textbf{Cơ sở lập kế hoạch.} Roadmap kỹ thuật gốc gồm 106 task, 216 ngày công cho phạm vi mở rộng, tính cho 7 người $\times$ 4 ngày hiệu dụng mỗi tuần. Nhóm có 5 thành viên (khoảng 20 ngày công mỗi tuần theo cùng giả định) và MVP chỉ gồm F1--F7, nên chỉ chọn 32 task của roadmap phục vụ F1--F7 (kể cả task tiên quyết), cộng 4 task mới chưa có trong roadmap (NEW-01 đến NEW-04). Tổng 32 task roadmap là 59,5 ngày công, cộng 4 task mới (NEW-01 đến NEW-04) là 9 ngày công, tổng cộng 68,5 ngày công. Mốc thời gian lấy theo tuần W0--W8 của roadmap (W0 bắt đầu 12/10/2026) và sẽ điều chỉnh theo hạn nộp Assignment 1: 20/10/2026.

\textbf{Quan hệ với phiên bản hiện có.} Một phần các chức năng F1--F3 và F7 đã có trong phiên bản DutyLog hiện tại (mục 6); các task bên dưới là phần cần kiểm kê, sửa hoặc bổ sung, nên trạng thái ``Review''/``Đang làm'' của một số task phản ánh mã đã viết sẵn. Task F-02 (kiểm kê API) là bước xác định phần nào còn thiếu.


\textbf{Task ngoài MVP (giữ ở Backlog).} Đăng ký, hủy và đổi ca (BE-01 đến BE-08, MB-03, MB-04); chế độ offline (BE-14, BE-15, MB-07); outbox, Kafka, thông báo đẩy và bảng trực tiếp (BE-16 đến BE-19, QR-04 đến QR-14); CI/CD nâng cao (CI-04 đến CI-16); kiểm thử tải và chaos (QA-02 đến QA-09). Một số hạng mục này đã có mã viết sẵn ở trạng thái Review (chưa chạy, chưa merge) nhưng không tính vào MVP cho đến khi có bằng chứng nhu cầu (mục 7.2).

\textbf{Phân công theo phần báo cáo.}

\medskip
{\centering
\begin{tabularx}{\dimexpr\textwidth-22pt\relax}{|>{\raggedright\arraybackslash}p{0.24\textwidth}|>{\centering\arraybackslash}p{0.12\textwidth}|X|}
\hline
\textbf{Thành viên} & \textbf{MSSV} & \textbf{Phụ trách (yêu cầu đề / mục báo cáo)} \\ \hline
Nguyễn Phúc Nhân & 2312438 & 2.1 Định hướng dự án (mục 1); 2.12 Nền tảng và kế hoạch dự án (mục 12) \\ \hline
Nguyễn Chu Hải Nam & 2312168 & 2.9 MVP và ưu tiên (mục 9); 2.10 User Stories (mục 10); 2.11 Luồng người dùng (mục 11) \\ \hline
Huỳnh Anh Chí Kiệt & 2013563 & 2.3 Giả định quan trọng (mục 3); 2.4 Kế hoạch bằng chứng (mục 4) \\ \hline
Nguyễn Văn Công Thành & 2313133 & 2.5 Nghiên cứu người dùng và thị trường (mục 5); 2.7 Đánh giá bằng chứng (mục 7) \\ \hline
Nguyễn Trần Trung Kiên & 2411759 & 2.2 Lý do chọn nền tảng di động (mục 2); 2.6 Giải pháp hiện có (mục 6); 2.8 Business Canvas (mục 8) \\ \hline
\end{tabularx}\par}

Phân công vai trò kỹ thuật trong các bảng M1--M4: Tech Lead: Nguyễn Phúc Nhân; Backend / Data: Huỳnh Anh Chí Kiệt; Mobile: Nguyễn Văn Công Thành; Product/UX: Nguyễn Chu Hải Nam; DevOps/SRE, QA/Perf \& Security: Nguyễn Trần Trung Kiên. Với task có hai vai trò (ví dụ ``Mobile + Backend''), vai trò đứng trước là người chịu trách nhiệm chính.

\textbf{Tải công việc theo vai trò (ngày công, tính cả task chung cho người chịu trách nhiệm chính):} Tech Lead $\approx$ 6 (kể cả NEW-04); Backend/Data $\approx$ 14,5 (kể cả QR-01); Mobile $\approx$ 21 (kể cả NEW-01 và NEW-03); Product/UX 14,5; DevOps/QA/Security 12,5. Mobile là vai trò nặng nhất và nằm trên đường găng của F3, nên Tech Lead nhận NEW-04 và hỗ trợ review/tích hợp API cho MB-01, NEW-01, NEW-03; tải được cập nhật trên GitHub Project ở đầu M2.


\bigskip
\textbf{M0 -- Hoàn thiện báo cáo Assignment 1} (19/10/2026): các mục 1--12 theo bảng phân công ở trên, kiểm tra liên kết GitHub trước khi nộp.

\bigskip
\textbf{M1 -- Khóa baseline và nền tảng an toàn (W0--W1, 12/10--25/10/2026).} Điều kiện hoàn thành: baseline và contract API được lưu trong Git; ghi nhận (check-in/check-out) yêu cầu JWT, kiểm tra quyền theo tổ chức/phòng, và không ghi trùng khi gửi lại.

{\footnotesize
\centering
\begin{tabularx}{\dimexpr\textwidth-22pt\relax}{|>{\raggedright\arraybackslash}p{0.09\textwidth}|X|>{\raggedright\arraybackslash}p{0.12\textwidth}|>{\raggedright\arraybackslash}p{0.14\textwidth}|>{\centering\arraybackslash}p{0.07\textwidth}|>{\raggedright\arraybackslash}p{0.16\textwidth}|}
\hline
\textbf{ID} & \textbf{Công việc} & \textbf{Chức năng} & \textbf{Phụ trách} & \textbf{Ngày} & \textbf{Trạng thái} \\ \hline
F-01 & Khóa baseline và tag trước khi sửa & Nền tảng & Tech Lead & 0,5 & Chưa làm \\ \hline
F-02 & Inventory: API thực tế so với tài liệu & Nền tảng & Tech Lead & 1 & Đang làm \\ \hline
F-03 & Đóng băng API contract v1 hiện tại & Nền tảng & Tech Lead & 1 & Chưa làm \\ \hline
F-04 & Định nghĩa FSM presence và shift & F2, F3, F6 & Tech Lead & 1 & Đang làm \\ \hline
F-05 & Phân định ownership dữ liệu và namespace & F1 & Tech Lead & 0,5 & Chưa làm \\ \hline
F-07 & Xác định kích thước tải thực & Nền tảng & Tech Lead & 1 & Chưa làm \\ \hline
CI-01 & Audit GitHub Actions ở hai repo & Chung & DevOps/SRE & 1 & Cần xác minh \\ \hline
CI-02 & Go lint, unit test, race, vet & Chung & DevOps/SRE & 1,5 & Chưa làm \\ \hline
CI-03 & Mobile typecheck, lint, unit test & Chung & DevOps/SRE & 1,5 & Đang làm \\ \hline
BE-09 & Bắt buộc JWT ở thao tác ghi & F1, F3 & Backend & 2 & Cần xác minh \\ \hline
BE-10 & Phân quyền theo tổ chức/phòng cho người trực & F1, F3 & Backend & 2 & Chưa làm \\ \hline
BE-12 & Idempotency cho check-in/check-out & F3, F4 & Backend & 2,5 & Chưa làm \\ \hline
BE-13 & Nguồn chân lý DB và nhất quán Redis & F7 & Backend & 2 & Cần xác minh \\ \hline
QR-01 & Truy vết đường quét: camera $\rightarrow$ API $\rightarrow$ DB $\rightarrow$ Redis & F3 & Backend + Data & 2 & Cần xác minh \\ \hline
QA-01 & Đo baseline máy quét trên thiết bị thật & F3 & QA/Perf & 2 & Chưa làm \\ \hline
UX-01 & Phỏng vấn thành viên và ca trưởng (mở rộng sang ban chủ nhiệm cho A7) & F7 & Product/UX & 2 & Chưa làm \\ \hline
\end{tabularx}\par
}

\bigskip
\textbf{M2 -- Thiết kế được kiểm chứng và luồng chính (W2--W3, 26/10--08/11/2026).} Điều kiện hoàn thành: người dùng thực hiện được đăng nhập, bắt đầu ca và quét một thành viên trên prototype hoặc bản chạy được; kết quả usability test được ghi lại.

{\footnotesize
\centering
\begin{tabularx}{\dimexpr\textwidth-22pt\relax}{|>{\raggedright\arraybackslash}p{0.09\textwidth}|X|>{\raggedright\arraybackslash}p{0.12\textwidth}|>{\raggedright\arraybackslash}p{0.14\textwidth}|>{\centering\arraybackslash}p{0.07\textwidth}|>{\raggedright\arraybackslash}p{0.16\textwidth}|}
\hline
\textbf{ID} & \textbf{Công việc} & \textbf{Chức năng} & \textbf{Phụ trách} & \textbf{Ngày} & \textbf{Trạng thái} \\ \hline
BE-11 & Bảo mật token QR và chống replay & F3, F5 & Backend & 2 & Cần xác minh \\ \hline
UX-02 & Quan sát ca trực tại 3 phòng (605 C6, B9, 710 H6) & F3, F7 & Product/UX & 2 & Chưa làm \\ \hline
UX-03 & Sitemap và điều hướng & F2, F7 & Product/UX & 1,5 & Review \\ \hline
UX-04 & Review prototype theo flow & F1--F7 & Product/UX & 2 & Đang làm \\ \hline
UX-12 & Design tokens và component library & Chung & Product/UX & 2,5 & Review \\ \hline
UX-05 & Usability test 5 người trên prototype & F3 & Product/UX & 2,5 & Chưa làm \\ \hline
UX-06 & Trạng thái loading/empty/error/permission & F3, F4 & Product/UX & 2 & Đang làm \\ \hline
MB-01 & Đăng nhập, giữ auth contract & F1 & Mobile & 2 & Review \\ \hline
MB-02 & Màn Trang chủ và điều hướng; hiển thị trạng thái phòng & F2, F7 & Mobile & 3 & Review \\ \hline
MB-05 & Màn quét và khung camera & F3, F4 & Mobile & 3 & Chưa làm \\ \hline
NEW-01 & Nhập MSSV thủ công, đánh dấu bản ghi ``nhập thủ công'' & F4 & Mobile + Backend & 3 & Chưa làm \\ \hline
NEW-04 & Nhập lịch ca ban đầu cho pilot (import CSV hoặc seed) và quy định người chịu trách nhiệm nhập lịch & F2 & Tech Lead & 1 & Chưa làm \\ \hline
\end{tabularx}\par
}

\bigskip
\textbf{Cổng quyết định G1 (cuối M2, 08/11/2026).} Dựa trên kết quả phỏng vấn (UX-01), quan sát ca trực (UX-02), thử nguyên mẫu (UX-05) và phân tích dữ liệu nhóm chat (mục 4.3), nhóm áp dụng tiêu chí ở mục 4.5 để kết luận A1. \emph{Support}: tiếp tục M3 đầy đủ. \emph{Challenge} (người dùng chỉ cần biết phòng mở/đóng và ai đang trực, không cần danh sách người có mặt): thu hẹp MVP về F1, F2 và F7 rút gọn (trạng thái mở/đóng theo ca của F2, người trực, lịch ca kế tiếp và thay đổi lịch; bỏ danh sách người có mặt) và hoãn luồng quét F3, F4, F5, F6. \emph{Unresolved}: sửa nguyên mẫu và thử lại trước khi đầu tư thêm cho F3. Các task F3 nặng (MB-05, MB-06, QA-01) vì vậy chỉ nên hoàn thiện sau G1; trước G1 chỉ làm phần chung với F4 và nền tảng.

\bigskip
\textbf{M3 -- Hoàn thiện MVP (W4--W6, 09/11--29/11/2026).} Điều kiện hoàn thành: đủ F1--F7 chạy được trên thiết bị thật; lỗi quét hoặc mất mạng hiển thị rõ, không báo thành công giả.

{\footnotesize
\centering
\begin{tabularx}{\dimexpr\textwidth-22pt\relax}{|>{\raggedright\arraybackslash}p{0.09\textwidth}|X|>{\raggedright\arraybackslash}p{0.12\textwidth}|>{\raggedright\arraybackslash}p{0.14\textwidth}|>{\centering\arraybackslash}p{0.07\textwidth}|>{\raggedright\arraybackslash}p{0.16\textwidth}|}
\hline
\textbf{ID} & \textbf{Công việc} & \textbf{Chức năng} & \textbf{Phụ trách} & \textbf{Ngày} & \textbf{Trạng thái} \\ \hline
MB-06 & Thông báo kết quả quét với trạng thái minh bạch & F3 & Mobile & 2 & Chưa làm \\ \hline
MB-08 & QR cá nhân: tự làm mới, đếm ngược, hết hạn & F5 & Mobile & 2 & Chưa làm \\ \hline
NEW-02 & Tự kết thúc ca và tự check-out khi hết giờ & F6 & Backend & 2 & Chưa làm \\ \hline
NEW-03 & Trạng thái phòng: ca kế tiếp và thay đổi lịch & F7 & Mobile + Backend & 3 & Chưa làm \\ \hline
MB-10 & Kiểm thử màn hình/e2e trên Android và iOS & Chung & Mobile & 3 & Chưa làm \\ \hline
QA-10 & Định nghĩa SLO và tiêu chí go/no-go & Chung & QA/Perf & 1,5 & Chưa làm \\ \hline
\end{tabularx}\par
}

\bigskip
\textbf{M4 -- Pilot và thu số liệu (W7--W8, 30/11--13/12/2026).} Điều kiện hoàn thành: chạy thử tại 3 phòng (605 C6 và B9 ở Cơ sở 1; 710 H6 ở Cơ sở 2) và thu các số đo của mục 4.5 (độ chính xác danh sách người có mặt, tỷ lệ ca có ghi nhận, thời gian check-in, tần suất mở trạng thái phòng, số lần nhập MSSV thủ công, số ca tự kết thúc).

{\footnotesize
\centering
\begin{tabularx}{\dimexpr\textwidth-22pt\relax}{|>{\raggedright\arraybackslash}p{0.09\textwidth}|X|>{\raggedright\arraybackslash}p{0.12\textwidth}|>{\raggedright\arraybackslash}p{0.14\textwidth}|>{\centering\arraybackslash}p{0.07\textwidth}|>{\raggedright\arraybackslash}p{0.16\textwidth}|}
\hline
\textbf{ID} & \textbf{Công việc} & \textbf{Chức năng} & \textbf{Phụ trách} & \textbf{Ngày} & \textbf{Trạng thái} \\ \hline
OPS-05 & Runbook sự cố và khôi phục & Chung & SRE + Security & 2 & Chưa làm \\ \hline
OPS-09 & Pilot 3 phòng và sign-off & F1--F7 & SRE + Security & 3 & Chưa làm \\ \hline
\end{tabularx}\par
}

\subsection{Bảng liên kết từ bằng chứng đến công việc}
Mỗi hàng đi từ bằng chứng (mục 7.1) đến task. Đầu vào của bảng là nhu cầu N1--N7 và yêu cầu R1--R6 (mục 10.1), story US1--US8 (mục 10.2), feature F1--F7 (mục 9.3) và hai flow (mục 11).

{\footnotesize
\centering
\begin{tabularx}{\dimexpr\textwidth-22pt\relax}{|>{\raggedright\arraybackslash}p{0.13\textwidth}|>{\raggedright\arraybackslash}p{0.06\textwidth}|>{\raggedright\arraybackslash}p{0.07\textwidth}|>{\raggedright\arraybackslash}p{0.12\textwidth}|>{\raggedright\arraybackslash}p{0.11\textwidth}|>{\centering\arraybackslash}p{0.06\textwidth}|X|}
\hline
\textbf{Nguồn thông tin} & \textbf{Nhu cầu} & \textbf{Yêu cầu} & \textbf{Câu chuyện người dùng} & \textbf{Chức năng} & \textbf{\scriptsize Luồng} & \textbf{Công việc} \\ \hline
Yêu cầu nền tảng (ngoài A1--A7) & --- & R1 & US1 & F1 & 1, 2 & F-05; BE-09; BE-10; MB-01 \\ \hline
A1 (7.1) & N1 & R6 & US7 & F7 & 2 & BE-13; MB-02; UX-01 \\ \hline
A2 (7.1) & N2 & R2, R5 & US2, US3, US6 & F2, F3, F6 & 1 & F-04; MB-05; MB-06; NEW-02 \\ \hline
A3 (7.1) & N3 & R3, R4 & US3, US4, US5 & F3, F4, F5 & 1 & BE-11; BE-12; MB-05; MB-08; NEW-01; QA-01 \\ \hline
A4, A5 (7.1) & N4, N7 & R2, R6 & US2, US7 & F2, F7 & 1, 2 & F-04; MB-02; NEW-03; NEW-04 \\ \hline
A6 (7.1) & N5 & R4, R5 & US3, US4, US5, US6 & F3, F4, F5, F6 & 1 & NEW-01; NEW-02; MB-06; UX-06; QA-01 \\ \hline
A7 (7.1, giả thuyết) & N6 & R6 & US8 & F7 & 2 & UX-01 (mở rộng sang ban chủ nhiệm); MB-02 \\ \hline
A1--A7 (kiểm chứng) & --- & --- & --- & F1--F7 & 1, 2 & F-01; F-02; F-03; F-07; CI-01; CI-02; CI-03; QR-01; UX-02; UX-03; UX-04; UX-05; UX-12; MB-10; QA-10; OPS-05; OPS-09 \\ \hline
\end{tabularx}\par
}